% Preâmbulo compartilhado pelos relatórios T02 - Redes de Dados
\documentclass[11pt,a4paper,oneside]{article}

\usepackage[utf8]{inputenc}
\usepackage[T1]{fontenc}
\usepackage[brazilian]{babel}
\usepackage{lmodern}
\usepackage{microtype}

\usepackage[a4paper,margin=2.5cm,top=2.5cm,bottom=2.5cm]{geometry}
\usepackage{graphicx}
\usepackage{xcolor}
\usepackage{tabularx}
\usepackage{longtable}
\usepackage{booktabs}
\usepackage{array}
\usepackage{enumitem}
\usepackage{titlesec}
\usepackage{fancyhdr}
\usepackage{listings}
\usepackage{hyperref}
\usepackage{caption}
\usepackage{float}
\usepackage{tcolorbox}
\tcbuselibrary{skins,breakable}

% Cores
\definecolor{codebg}{RGB}{245,245,247}
\definecolor{codeborder}{RGB}{220,220,225}
\definecolor{codekey}{RGB}{0,90,160}
\definecolor{codecom}{RGB}{100,110,120}
\definecolor{codestr}{RGB}{40,140,80}
\definecolor{noteblue}{RGB}{30,80,160}
\definecolor{notebg}{RGB}{235,243,255}
\definecolor{warnyellow}{RGB}{170,120,0}
\definecolor{warnbg}{RGB}{255,248,225}

% Hyperref
\hypersetup{
  colorlinks=true,
  linkcolor=black,
  urlcolor=codekey,
  citecolor=black,
  pdfborder={0 0 0}
}

% Cabeçalho/rodapé
\pagestyle{fancy}
\fancyhf{}
\fancyhead[L]{\small T02 - Redes de Dados}
\fancyhead[R]{\small \leftmark}
\fancyfoot[C]{\thepage}
\renewcommand{\headrulewidth}{0.4pt}

% Listings (código)
\lstdefinestyle{shell}{
  backgroundcolor=\color{codebg},
  basicstyle=\ttfamily\footnotesize,
  keywordstyle=\color{codekey}\bfseries,
  commentstyle=\color{codecom}\itshape,
  stringstyle=\color{codestr},
  numbers=none,
  frame=single,
  rulecolor=\color{codeborder},
  framesep=4pt,
  xleftmargin=6pt,
  xrightmargin=6pt,
  breaklines=true,
  breakatwhitespace=false,
  showstringspaces=false,
  columns=fullflexible,
  keepspaces=true,
  literate=
    {á}{{\'a}}1 {é}{{\'e}}1 {í}{{\'i}}1 {ó}{{\'o}}1 {ú}{{\'u}}1
    {Á}{{\'A}}1 {É}{{\'E}}1 {Í}{{\'I}}1 {Ó}{{\'O}}1 {Ú}{{\'U}}1
    {ã}{{\~a}}1 {õ}{{\~o}}1 {Ã}{{\~A}}1 {Õ}{{\~O}}1
    {â}{{\^a}}1 {ê}{{\^e}}1 {ô}{{\^o}}1 {Â}{{\^A}}1 {Ê}{{\^E}}1 {Ô}{{\^O}}1
    {ç}{{\c{c}}}1 {Ç}{{\c{C}}}1
    {à}{{\`a}}1
}
\lstset{style=shell}

% Caixas de nota/observação/atenção
\newtcolorbox{notebox}{
  enhanced, breakable,
  colback=notebg, colframe=noteblue,
  boxrule=0.6pt, arc=2pt,
  left=8pt, right=8pt, top=4pt, bottom=4pt,
  fonttitle=\bfseries
}
\newtcolorbox{warnbox}{
  enhanced, breakable,
  colback=warnbg, colframe=warnyellow,
  boxrule=0.6pt, arc=2pt,
  left=8pt, right=8pt, top=4pt, bottom=4pt,
  fonttitle=\bfseries
}

% Imagem local
% Uso: \imgph{largura}{caminho}
\newcommand{\imgph}[2]{%
  \begin{figure}[H]\centering
    \includegraphics[width=#1]{#2}%
  \end{figure}
}

% Espaçamento de seções
\titleformat{\section}{\Large\bfseries\color{noteblue}}{\thesection.}{0.6em}{}
\titleformat{\subsection}{\large\bfseries}{\thesubsection}{0.6em}{}
\titleformat{\subsubsection}{\normalsize\bfseries}{\thesubsubsection}{0.6em}{}

% Capa
\newcommand{\capa}[2]{%
  \begin{titlepage}
    \centering
    \vspace*{1.5cm}
    % Logo: coloque o arquivo em .tex/img/logo.png
    \IfFileExists{imgs/logo.png}{\includegraphics[width=4cm]{imgs/logo.png}}{%
      \fbox{\parbox[c][3cm][c]{4cm}{\centering\footnotesize\itshape [logo.png]}}}
    \vfill
    {\large\bfseries T02 -- Redes de Dados\par}
    \vspace{1cm}
    {\Huge\bfseries\color{noteblue} #1\par}
    \vspace{0.8cm}
    {\large #2\par}
    \vfill
    {\large \today\par}
    \vspace{1cm}
  \end{titlepage}
}


\begin{document}

\capa{Relatório 11}{Configuração de IDS de Rede usando Suricata no Debian (VM no VirtualBox)}

\tableofcontents
\newpage

\noindent
Este relatório descreve a configuração de um \textbf{IDS} com \textbf{Suricata} \textbf{diretamente no Debian (uma única VM)}, monitorando o tráfego da própria VM para a \textbf{Internet}. Usaremos \textbf{apenas regras locais (Custom Rules)}, sem baixar rulesets externos. O objetivo é \textbf{instalar}, \textbf{iniciar} o Suricata em \textbf{modo IDS}, \textbf{gerar tráfego didático} (via navegador/\texttt{curl}) e \textbf{analisar os alertas} no \texttt{fast.log}/\texttt{eve.json}.

\begin{itemize}
  \item \textbf{Material de apoio:} \url{https://github.com/vin1sss/T02---Redes-de-Dados-1/}
\end{itemize}

\section{Introdução}

A única VM (Debian Desktop) atuará como \textbf{sensor IDS} (Suricata) e como \textbf{cliente} que gera tráfego. As regras locais foram desenhadas para disparar alertas de forma controlada e didática (URI, BODY/POST, HEADER e DNS).

Um IDS de rede opera de forma passiva: observa os pacotes, aplica assinaturas/regras de inspeção e gera alertas quando encontra padrões suspeitos, sem bloquear diretamente o tráfego. No Suricata, essa detecção ocorre com inspeção profunda de protocolos (DPI), permitindo visibilidade de elementos como URI, cabeçalhos, corpo HTTP e consultas DNS.

\textbf{Pilares:} \textbf{Confidencialidade} e \textbf{Monitoramento/Detecção}, com ênfase em \textbf{auditoria} e \textbf{visibilidade} do tráfego.

\section{Conceitos e Fundamentos}

\begin{itemize}
  \item \textbf{Firewall $\times$ IDS:} firewall \textbf{impõe} política (permite/bloqueia); IDS \textbf{observa e alerta} (passivo).
  \item \textbf{Suricata (DPI):} inspeção por \textbf{regras/assinaturas}. Vamos criar \textbf{regras locais} para comprovar a capacidade de inspeção de \textbf{URI}, \textbf{cabeçalhos}, \textbf{corpo HTTP} e \textbf{consultas DNS}.
  \item \textbf{HTTPS vs. HTTP:} em \textbf{HTTPS} o conteúdo é cifrado (IDS vê metadados, p.ex. SNI/JA3); em \textbf{HTTP} o IDS consegue inspecionar cabeçalhos e corpo.
  \item \textbf{Termos das regras (resumo):}
  \begin{itemize}
    \item \texttt{flow:to\_server} = tráfego \textbf{do cliente para o servidor} (requisição).
    \item \texttt{http.uri} = caminho solicitado (p.ex. \texttt{/rota}).
    \item \texttt{http\_header} = cabeçalhos HTTP (ex.: \texttt{User-Agent}, \texttt{Host}, cabeçalhos customizados).
    \item \texttt{http.request\_body} = corpo enviado pelo cliente (POST/PUT).
    \item \texttt{dns.query} = nome de domínio consultado no DNS.
    \item \texttt{classtype}/\texttt{priority} = \textbf{classificação} e \textbf{prioridade} do alerta no log.
  \end{itemize}
\end{itemize}

\section{Ambiente e Topologia}

\textbf{Topologia lógica (1 VM):}
\begin{itemize}
  \item \textbf{user 1 - Debian (com Suricata)} atuando como sensor IDS e cliente de teste.
  \item Saída para Internet através de NAT.
\end{itemize}

\begin{notebox}
\textbf{Topologia:} \textbf{user 1 - Debian $\leftrightarrow$ NAT (VBox) $\leftrightarrow$ Internet}.
\end{notebox}

\subsection{Resumo do laboratório}

\begin{tabularx}{\linewidth}{lX}
\toprule
\textbf{Elemento} & \textbf{Definição} \\
\midrule
Rede do VirtualBox & \textbf{NAT} com DHCP do VirtualBox \\
VM do laboratório & \textbf{user 1 - Debian} \\
Papel da VM & Sensor IDS com \textbf{Suricata} e cliente que gera o tráfego \\
Tráfego observado & Saída da própria VM para a Internet \\
Evidências principais & Alertas em \texttt{fast.log} e eventos em \texttt{eve.json} \\
Recursos da VM & \textbf{6 GB de RAM} e \textbf{3 núcleos de CPU} \\
\bottomrule
\end{tabularx}

\subsection{Referências usadas durante a execução}

\begin{tabularx}{\linewidth}{lX}
\toprule
\textbf{Item} & \textbf{Uso} \\
\midrule
Senha das máquinas virtuais & \texttt{pytest} \\
Senha solicitada pelo \texttt{sudo} & \texttt{pytest} \\
Interface de saída da VM & Confirmada com \texttt{ip route get 1.1.1.1} \\
\texttt{enp0s3} & Nome usado nos exemplos do relatório para a interface monitorada \\
\texttt{/etc/suricata/rules/local.rules} & Arquivo das regras locais do experimento \\
\texttt{/var/log/suricata/fast.log} & Log textual acompanhado durante os testes \\
\texttt{/var/log/suricata/eve.json} & Log JSON usado para consulta alternativa dos alertas \\
\bottomrule
\end{tabularx}

\subsection{Pacotes usados no laboratório}

\begin{itemize}
  \item \texttt{dnsutils}, \texttt{net-tools}, \texttt{curl}, \texttt{wget}
  \item \texttt{suricata}, \texttt{tcpdump}, \texttt{jq}
\end{itemize}

\subsection{Guia auxiliar de VirtualBox}

\begin{notebox}
Se não estiver conseguindo copiar e colar comandos entre o computador host e a VM, redimensionar a tela corretamente ou usar melhor a integração com o VirtualBox, consulte o guia \texttt{VirtualBox\_Guest\_Additions.md} para instalar o \textbf{VirtualBox Guest Additions}.
\end{notebox}

\section{Instalação e Preparação}

\begin{notebox}
Se a interface indicada por \texttt{ip route get 1.1.1.1} não for \texttt{enp0s3}, substitua apenas esse nome nos comandos de inicialização e captura que usam a interface de rede. Mantenha essa interface anotada antes de seguir.
\end{notebox}

\subsection{Preparar o ambiente no VirtualBox}

\begin{itemize}
  \item \textbf{VM:} \texttt{user 1 - Debian}
  \item \textbf{Adapter 1:} \textbf{NAT} (DHCP do VBox)
\end{itemize}

\imgph{0.65\linewidth}{https://github.com/user-attachments/assets/518f2ddb-4262-4f51-bb8d-7c9dd8cb2e6b}

\begin{itemize}
  \item \textbf{Recursos da VM:} \textbf{6 GB de RAM} e \textbf{3 núcleos de CPU}
\end{itemize}

\imgph{0.4\linewidth}{https://github.com/user-attachments/assets/11c305c2-950f-4493-996b-427af0095074}

\begin{itemize}
  \item \textbf{Iniciar a VM}
\end{itemize}

\begin{notebox}
\textbf{Ao final desta etapa:} a VM deve estar configurada em \textbf{NAT}, com os recursos previstos e pronta para receber endereçamento via DHCP.
\end{notebox}

\subsection{Após iniciar a VM, conferir a rede e padronizar o DHCP}

\subsubsection*{A) Padronizar a interface em DHCP}

\begin{lstlisting}
IFACE=$(ip route | awk '/default/ {print $5; exit}')
sudo dhclient -r "$IFACE" || true
sudo ip addr flush dev "$IFACE"
sudo dhclient "$IFACE"
ip -4 a show "$IFACE"
\end{lstlisting}

\imgph{0.6\linewidth}{https://github.com/user-attachments/assets/905455e8-4a02-4a6a-ab2e-5e7f0ebad5d7}

\subsubsection*{B) Validar IP, saída para Internet e interface monitorada}

\begin{lstlisting}
ip a                    # deve obter IP 10.0.2.x (NAT do VBox)
ping -c2 1.1.1.1        # checar saida para internet
ip route get 1.1.1.1    # confirmar interface de saida (ex.: enp0s3)
\end{lstlisting}

\imgph{0.7\linewidth}{https://github.com/user-attachments/assets/50ad173d-5d56-426f-bc52-58c3e6b848a3}

\begin{notebox}
\textbf{Resultado esperado:} a VM apresenta IP coerente com a rede NAT, alcança \texttt{1.1.1.1} e indica a interface de saída usada pelo tráfego observado.

\textbf{Ao final desta etapa:} a conectividade externa e a interface de monitoramento devem estar confirmadas antes da instalação do Suricata.
\end{notebox}

\subsection{Preparar o sistema antes da instalação}

\subsubsection*{A) Prevenir lock do \texttt{apt}}

\begin{lstlisting}
sudo systemctl stop packagekit || true
sudo pkill -x packagekitd || true
sudo dpkg --configure -a
\end{lstlisting}

\imgph{0.45\linewidth}{https://github.com/user-attachments/assets/d30e0ae4-9ecf-491a-8501-b15edcc568c7}

\begin{notebox}
\textbf{Ao final desta etapa:} o sistema deve estar liberado para instalar os pacotes sem conflito de gerenciador.
\end{notebox}

\subsection{Instalar utilitários e Suricata}

\begin{lstlisting}
sudo apt update && sudo apt install -y dnsutils net-tools curl wget suricata tcpdump jq
\end{lstlisting}

\imgph{0.75\linewidth}{https://github.com/user-attachments/assets/4ea2ec37-be75-4134-965c-24403e454204}

\begin{notebox}
\textbf{Ao final desta etapa:} todas as ferramentas para gerar tráfego, monitorar a rede e consultar alertas devem estar instaladas.
\end{notebox}

\subsection{Criar e validar as regras locais}

\subsubsection*{A) Entender o papel das regras}

\begin{notebox}
\textbf{O que cada regra faz:}
\begin{itemize}
  \item \textbf{SID 1002001 (URI):} alerta quando a \textbf{URI} contém \texttt{/SURICATA\_BROWSER\_URI\_01}. Inspeção de caminho (HTTP em claro).
  \item \textbf{SID 1002002 (BODY/POST):} alerta quando o \textbf{corpo} contém \texttt{SURICATA\_BROWSER\_BODY\_02}. Inspeção de payload em POST/PUT.
  \item \textbf{SID 1002003 (HEADER):} alerta quando um \textbf{cabeçalho HTTP} contém \texttt{X-Trigger-Lab: 1}. Inspeção de cabeçalhos.
  \item \textbf{SID 1002004 (DNS):} alerta quando a \textbf{consulta DNS} contém \texttt{suricata-trigger-lab.example}. Visibilidade de DNS.
\end{itemize}
\textbf{Arquivo:} \texttt{/etc/suricata/rules/local.rules} (uma linha por regra).
\end{notebox}

\subsubsection*{B) Criar o arquivo \texttt{local.rules}}

\begin{lstlisting}
sudo tee /etc/suricata/rules/local.rules >/dev/null <<'RULES'
alert http any any -> any any (msg:"CUSTOM BROWSER - HTTP URI trigger"; flow:to_server; http.uri; content:"/SURICATA_BROWSER_URI_01"; nocase; sid:1002001; rev:4; classtype:web-application-activity; priority:2;)
alert http any any -> any any (msg:"CUSTOM BROWSER - HTTP client body trigger"; flow:to_server; http.request_body; content:"SURICATA_BROWSER_BODY_02"; nocase; sid:1002002; rev:4; classtype:web-application-attack; priority:2;)
alert http any any -> any any (msg:"CUSTOM BROWSER - HTTP header X-Trigger-Lab"; flow:to_server; content:"X-Trigger-Lab|3a| 1"; http_header; nocase; sid:1002003; rev:5; classtype:attempted-recon; priority:2;)
alert udp any any -> any 53 (msg:"CUSTOM BROWSER - DNS query trigger"; dns.query; content:"suricata-trigger-lab.example"; nocase; sid:1002004; rev:3; classtype:attempted-recon; priority:3;)
RULES
\end{lstlisting}

\imgph{0.85\linewidth}{https://github.com/user-attachments/assets/5ab11d4f-11c8-4de4-ac55-a93533803b58}

\subsubsection*{C) Validar a sintaxe das regras}

\begin{lstlisting}
sudo suricata -T -S /etc/suricata/rules/local.rules -v
# esperado: "4 rules successfully loaded, 0 failed"
\end{lstlisting}

\imgph{0.85\linewidth}{https://github.com/user-attachments/assets/480cb59f-ded9-49f3-bfbf-3d3b8d5fdedb}

\begin{notebox}
\textbf{Ao final desta etapa:} o arquivo de regras deve existir e o Suricata deve confirmar o carregamento das \textbf{4 regras locais}.
\end{notebox}

\subsection{Iniciar o Suricata em modo IDS e verificar a operação}

\begin{warnbox}
Em Debian, usar \texttt{--set outputs.*} pode causar erro de ``child node (null)''. Usaremos o YAML padrão.
\end{warnbox}

\subsubsection*{A) Subir o Suricata manualmente}

\begin{lstlisting}
# parar instancias anteriores e limpar pid/log
sudo systemctl stop suricata
sudo pkill -x suricata || true
sudo rm -f /var/run/suricata.pid /var/log/suricata/fast.log

# subir em daemon na interface correta (ex.: enp0s3)
sudo suricata -i enp0s3 -S /etc/suricata/rules/local.rules -l /var/log/suricata -D

# conferir que esta ativo
ps aux | grep '[s]uricata'
sudo tail -n 15 /var/log/suricata/suricata.log
\end{lstlisting}

\imgph{0.85\linewidth}{https://github.com/user-attachments/assets/bb721063-caf0-46e3-947c-1a2e5157cdea}

\subsubsection*{B) Confirmar onde os alertas serão acompanhados}

\textbf{Onde ver alertas:} \texttt{tail -f /var/log/suricata/fast.log}

\textbf{Alternativa JSON:} \texttt{jq 'select(.event\_type=="alert")' /var/log/suricata/eve.json | tail -n 5}

\begin{notebox}
\textbf{Ao final desta etapa:} o Suricata deve estar ativo em modo IDS, carregando \texttt{local.rules} e com logs prontos para a observação dos testes.
\end{notebox}

\section{Procedimentos (Passo a Passo --- user 1 - Debian $\leftrightarrow$ Internet)}

\subsection*{Checklist antes dos testes}

\begin{itemize}
  \item as \textbf{4 regras locais} foram validadas com \texttt{suricata -T};
  \item o processo do \textbf{Suricata} está em execução;
  \item o terminal de acompanhamento do \texttt{fast.log} está aberto;
  \item a interface monitorada foi identificada (ex.: \texttt{enp0s3});
  \item a VM possui conectividade para gerar o tráfego HTTP e DNS.
\end{itemize}

\subsection*{Organização dos terminais}

\textbf{Terminal A:} \texttt{tail -f /var/log/suricata/fast.log}

\imgph{0.5\linewidth}{https://github.com/user-attachments/assets/f7e2bad4-8a20-4919-979f-dd084d40962b}

\textbf{Terminal B} (caso optar por testar usando \texttt{curl} diretamente no terminal): comandos de disparo abaixo. Dica: force IPv4 com \texttt{-4}.

\subsection{Cenário 1 --- HTTP URI (SID 1002001)}

\textbf{O que este teste demonstra:} que o Suricata consegue \textbf{inspecionar a URI} (caminho) de uma requisição HTTP e gerar alerta quando encontra a \textbf{string-alvo}.

\subsubsection*{A) Executar o gatilho}

\begin{lstlisting}
curl -4 -s 'http://neverssl.com/SURICATA_BROWSER_URI_01' >/dev/null
# ou navegador: http://neverssl.com/SURICATA_BROWSER_URI_01
\end{lstlisting}

\begin{itemize}
  \item Caso demore para executar ou alertar, aguarde. Este comportamento é esperado.
\end{itemize}

\imgph{0.95\linewidth}{https://github.com/user-attachments/assets/75343aec-013b-4504-be75-92acb3f86fa6}

\subsubsection*{B) Observar o alerta no log}

\textbf{O que observar:} uma linha no \texttt{fast.log} com a mensagem \texttt{CUSTOM BROWSER - HTTP URI trigger} (SID 1002001), indicando tráfego \texttt{\{TCP\} <IP\_VM>:<porta> -> <IP\_destino>:80}.

\begin{notebox}
\textbf{Ao final deste cenário:} deve existir evidência de que a regra de \textbf{URI} foi acionada pelo tráfego HTTP gerado.
\end{notebox}

\subsection{Cenário 2 --- HTTP BODY/POST (SID 1002002)}

\textbf{O que este teste demonstra:} que o Suricata consegue \textbf{inspecionar o corpo} (payload) de uma requisição HTTP (POST/PUT) e alertar quando a \textbf{string} estiver presente.

\subsubsection*{A) Executar o gatilho}

\begin{lstlisting}
curl -4 -s -X POST -d 'SURICATA_BROWSER_BODY_02' 'http://neverssl.com/' >/dev/null
# (opcional para ver o request)
# curl -4 -v -X POST -d 'SURICATA_BROWSER_BODY_02' 'http://neverssl.com/' >/dev/null
\end{lstlisting}

\imgph{0.95\linewidth}{https://github.com/user-attachments/assets/24087409-8c8b-4a65-b4f8-8fe91cf8ccb7}

\subsubsection*{B) Observar o alerta no log}

\textbf{O que observar:} a mensagem \texttt{CUSTOM BROWSER - HTTP client body trigger} (SID 1002002). Se não disparar de primeira, rode com \texttt{-v} e valide que o \textbf{POST} saiu; rode novamente.

\begin{notebox}
\textbf{Ao final deste cenário:} deve existir evidência de que a regra de \textbf{corpo da requisição} foi acionada.
\end{notebox}

\subsection{Cenário 3 --- HTTP HEADER (SID 1002003)}

\textbf{O que este teste demonstra:} que o Suricata consegue \textbf{inspecionar cabeçalhos HTTP} e gerar alerta quando um \textbf{cabeçalho customizado} (\texttt{X-Trigger-Lab: 1}) está presente.

\subsubsection*{A) Executar o gatilho}

\begin{lstlisting}
curl -4 -s -H 'X-Trigger-Lab: 1' 'http://neverssl.com/' >/dev/null
# (opcional verbose)
# curl -4 -v -H 'X-Trigger-Lab: 1' 'http://neverssl.com/' >/dev/null
\end{lstlisting}

\imgph{0.95\linewidth}{https://github.com/user-attachments/assets/6f87fb49-cc08-4880-9b18-6cc2c3657e72}

\subsubsection*{B) Observar o alerta no log}

\textbf{O que observar:} a mensagem \texttt{CUSTOM BROWSER - HTTP header X-Trigger-Lab} (SID 1002003), com fluxo \texttt{\{TCP\} <IP\_VM>:<porta> -> <IP\_destino>:80}.

\subsection{Cenário 4 --- DNS (SID 1002004)}

\textbf{O que este teste demonstra:} que o Suricata enxerga \textbf{consultas DNS} e alerta quando a consulta contém o \textbf{domínio-alvo}.

\subsubsection*{A) Executar o gatilho}

\begin{lstlisting}
dig +short suricata-trigger-lab.example >/dev/null
# ou navegador: http://suricata-trigger-lab.example/
\end{lstlisting}

\imgph{0.95\linewidth}{https://github.com/user-attachments/assets/5cb962e1-b239-4938-ac74-6255a2d782f4}

\subsubsection*{B) Observar o alerta no log}

\textbf{O que observar:} a mensagem \texttt{CUSTOM BROWSER - DNS query trigger} (SID 1002004), com fluxo \texttt{\{UDP\} <IP\_VM>:<porta> -> <DNS\_resolvedor>:53}.

\section{Verificação e Resultados}

\subsection{Checklist de sucesso do experimento}

\begin{itemize}
  \item \texttt{suricata -T} retorna \texttt{4 rules successfully loaded, 0 rules failed}.
  \item O Suricata permanece em execução (\texttt{ps aux | grep '[s]uricata'}) e sem erro crítico no \texttt{suricata.log}.
  \item Cada um dos quatro cenários gera o alerta correspondente em \texttt{fast.log} ou \texttt{eve.json}.
\end{itemize}

\subsection{Quadro comparativo (cenários)}

\begin{tabularx}{\linewidth}{cXXccc}
\toprule
\textbf{\#} & \textbf{Tráfego} & \textbf{Regra (msg)} & \textbf{SID} & \textbf{Resultado} & \textbf{Verificar em} \\
\midrule
1 & user 1 - Debian $\rightarrow$ Internet (HTTP/URI)    & \texttt{CUSTOM BROWSER - HTTP URI trigger}          & 1002001 & Alertas & \texttt{fast.log} / \texttt{eve.json} \\
2 & user 1 - Debian $\rightarrow$ Internet (HTTP/Body)   & \texttt{CUSTOM BROWSER - HTTP client body trigger}  & 1002002 & Alertas & \texttt{fast.log} / \texttt{eve.json} \\
3 & user 1 - Debian $\rightarrow$ Internet (HTTP/Header) & \texttt{CUSTOM BROWSER - HTTP header X-Trigger-Lab} & 1002003 & Alertas & \texttt{fast.log} / \texttt{eve.json} \\
4 & user 1 - Debian $\rightarrow$ Internet (DNS)         & \texttt{CUSTOM BROWSER - DNS query trigger}         & 1002004 & Alertas & \texttt{fast.log} / \texttt{eve.json} \\
\bottomrule
\end{tabularx}

\subsection{Coleta rápida dos alertas}

\begin{lstlisting}
# Ultimos alertas (formato fast)
tail -n 20 /var/log/suricata/fast.log

# Em JSON (EVE) com campos uteis, 1 por linha
jq -r 'select(.event_type=="alert") | "\(.timestamp) SID=\(.alert.signature_id) MSG=\(.alert.signature) \(.src_ip):\(.src_port) -> \(.dest_ip):\(.dest_port)"' \
  /var/log/suricata/eve.json | tail -n 10
\end{lstlisting}

\imgph{0.95\linewidth}{https://github.com/user-attachments/assets/8cbeaf01-02c1-49a8-9517-e756fd0f9a29}

\section{Conclusão}

Com \textbf{Suricata na própria VM} e \textbf{regras locais} simples, comprovamos a detecção de \textbf{URI}, \textbf{cabeçalhos}, \textbf{corpo HTTP} e \textbf{consultas DNS} por meio de cenários objetivos e reproduzíveis. A configuração evita dependências externas e foca em \textbf{didática}, \textbf{auditoria} e \textbf{visibilidade} do tráfego \textbf{HTTP/DNS}.

\section*{Apêndice --- Troubleshooting rápido}
\addcontentsline{toc}{section}{Apêndice --- Troubleshooting rápido}

\begin{itemize}
  \item \textbf{Sem alertas no \texttt{fast.log}:}
    \begin{enumerate}
      \item Confirme que o Suricata está \textbf{rodando}: \texttt{ps aux | grep '[s]uricata'}
      \item Veja o log de inicialização: \texttt{sudo tail -n 30 /var/log/suricata/suricata.log}
      \item Verifique \textbf{tráfego} na interface: \texttt{sudo tcpdump -ni enp0s3 'port 80 or port 53' -c 10}
    \end{enumerate}

  \item \textbf{Erro ``Failed to lookup configuration child node: (null)'' ao iniciar:} não use \texttt{--set outputs.*} no Debian. Inicie assim: \texttt{sudo suricata -i enp0s3 -S /etc/suricata/rules/local.rules -l /var/log/suricata -D}

  \item \textbf{\texttt{pidfile} ``stale'' / conflito de instância:}
\begin{lstlisting}
sudo systemctl stop suricata
sudo pkill -x suricata || true
sudo rm -f /var/run/suricata.pid
\end{lstlisting}

  \item \textbf{Regras não carregam / erros de sintaxe:} cada regra deve estar \textbf{em uma única linha} (Snort-syntax). Valide com \texttt{-T}.

  \item \textbf{HTTP não dispara (HSTS/HTTPS forçado):} use \texttt{http://neverssl.com} e force IPv4 com \texttt{curl -4}.

  \item \textbf{Confirmar que os triggers realmente estão saindo:} para POST/HEADER, rode em paralelo \texttt{sudo tcpdump -A -ni enp0s3 'port 80' -c 10} e verifique o \texttt{POST}, o corpo e o cabeçalho \texttt{X-Trigger-Lab} em texto.

  \item \textbf{Classificação diferente no log:} \texttt{classtype:attempted-recon} pode aparecer como \textbf{Attempted Information Leak} no \texttt{fast.log}.
\end{itemize}

\end{document}
